\chapter{Numerical reasoning: approximation with an error budget}
\label{learn:numerical}

\begin{learningcard}{Prerequisites and task}
Use Taylor bounds from Chapter~\ref{learn:calculus}, matrix maps from
Chapter~\ref{learn:linear_algebra}, and stability from
Chapter~\ref{learn:dynamics}. A numerical answer is useful when its precision
matches the question and its errors are understood. We ask how to distinguish
an accurate-looking output from an approximation that survives refinement
and independent checks.
\end{learningcard}

\section{Three errors can coexist}
Modeling error measures the difference between a chosen mathematical model
and the target phenomenon. Discretization or truncation error comes from
replacing an infinite or continuous operation by a finite one. Roundoff error
comes from representing and manipulating numbers with finite precision.
Changing the mesh can reduce discretization error while leaving a wrong model
unchanged. Exact rational arithmetic can remove roundoff for some calculations
while leaving truncation error intact.

A problem is well-conditioned when small perturbations of its input produce
small relative changes in the requested output, with scales specified.
For a differentiable scalar function away from zero input and output, the
relative condition number is $|xf'(x)/f(x)|$. An algorithm's stability
instead concerns how its errors behave while solving the problem. A stable
algorithm can still produce an uncertain answer for an ill-conditioned input.

\section{Worked example: a smaller difference can lose information}
Approximate $f'(x)$ by
$D_hf(x)=[f(x+h)-f(x)]/h$ for $h\ne0$. Taylor's theorem gives
$D_hf(x)=f'(x)+hf''(\xi)/2$ for a suitable intermediate point when the
required smoothness holds. If $|f''|\leq M$, truncation error is at most
$M|h|/2$.

Suppose each function evaluation has absolute error at most $\eta$.
Subtracting the evaluations and dividing by $h$ adds an error at most
$2\eta/|h|$, before accounting for the arithmetic operations themselves.
The combined illustrative bound is
\[
 |D_h^{\mathrm{computed}}f(x)-f'(x)|
 \leq \frac{M|h|}{2}+\frac{2\eta}{|h|}.
\]
For $M,\eta>0$, balancing these two terms suggests
$|h|=2\sqrt{\eta/M}$. Arbitrarily shrinking $h$ amplifies evaluation
error. This formula depends on the absolute-error model; a real program
must account for scales and the actual evaluation routine.

For $f(x)=x^2$ at $x=1$, exact algebra gives $D_hf(1)=2+h$.
At $h=0.01$, the error is exactly $0.01$. A floating-point implementation
may eventually round $1+h$ to one, producing a zero numerator despite the
true derivative two. An exact symbolic identity provides a better oracle
than merely comparing two nearby floating-point outputs.

\section{Worked example: integration with a computable remainder}
The composite trapezoidal rule on $[a,b]$ with $n$ equal pieces of width
$h=(b-a)/n$ is
\[
 T_n=h\left(\frac{f(a)+f(b)}2+\sum_{j=1}^{n-1}f(a+jh)\right).
\]
For $f$ with continuous second derivative bounded by $M$, its error obeys
$|T_n-\int_a^bf(x)\,dx|\leq(b-a)Mh^2/12$.
One obtains the order by integrating the local linear-interpolation error
on each piece: a local error of order $h^3$ accumulated across $n$ pieces
gives a global error of order $h^2$.

For $f(x)=x^2$ on $[0,1]$, use
$\sum_{j=1}^{n-1}j^2=(n-1)n(2n-1)/6$ to derive
\[
 T_n=\frac1{2n}+\frac{(n-1)n(2n-1)}{6n^3}
 =\frac13+\frac1{6n^2}.
\]
The exact integral is $1/3$, so the error is positive and equals
$1/(6n^2)$. Since $M=2$, the general bound is attained. Doubling $n$
reduces the error by four. The extrapolated quantity
$(4T_{2n}-T_n)/3$ equals $1/3$ exactly for this quadratic because the
leading error formula is exact. For other smooth functions, extrapolation
requires the appropriate error expansion.

\begin{learningcard}{Historical situation: reliable matrix computation}
Turing's 1948 paper ``Rounding-off Errors in Matrix Processes'' studies
error in numerical matrix operations. The paper's subject is a concrete
reminder that computation requires analysis of the procedure as well as
algebraic correctness of the equations. Its publication documents an early
analysis of such errors; it supplies neither the origin story of all
numerical analysis nor a guarantee for a modern implementation.
\footnote{Alan M. Turing, ``Rounding-off Errors in Matrix Processes,''
\emph{Quarterly Journal of Mechanics and Applied Mathematics} 1 (1948),
pp.~287--308: \url{https://doi.org/10.1093/qjmam/1.1.287}.}
The finite-difference experiment above uses a modern explicit error budget.
\end{learningcard}

\section{Iteration: a residual and an error answer different questions}
To solve $Ax=b$, consider $x_{m+1}=Bx_m+c$. If a fixed point $x_*$
exists, its error obeys $e_{m+1}=Be_m$. Any induced norm with $\|B\|<1$
gives $\|e_m\|\leq\|B\|^m\|e_0\|$, proving convergence. In finite
dimensions, spectral radius below one is the general constant-matrix criterion
for convergence from every initial error, although the chosen norm can show
transient growth before decay.

For the system $2x+y=3$, $x+2y=3$, simultaneous updates give
$x_{m+1}=(3-y_m)/2$, $y_{m+1}=(3-x_m)/2$. The error matrix is
\[
 B=\begin{pmatrix}0&-1/2\\-1/2&0\end{pmatrix}.
\]
Its Euclidean operator norm is $1/2$. Starting at $(0,0)$ produces
$(1.5,1.5)$, $(0.75,0.75)$, $(1.125,1.125)$, approaching $(1,1)$.
The residual is $r=b-A\widehat x$. If $A$ is invertible, the actual
error satisfies $\widehat x-x_*=-A^{-1}r$, so
$\|\widehat x-x_*\|\leq\|A^{-1}\|\|r\|$. A small residual establishes
a small solution error only when the inverse scale is controlled.

\section{A convergence study needs an oracle and a regime}
For a time-dependent scheme, consistency asks whether the discrete equation
approximates the differential equation as the spacings shrink. Stability
controls growth of perturbations. Convergence asks whether the computed
solution approaches the desired exact solution. The familiar equivalence
between consistency plus stability and convergence has specified hypotheses
for linear well-posed problems; it is not a universal statement about every
nonlinear simulation.

Conservation gives another independent test. As
Chapter~\ref{learn:dynamics} shows, the periodic explicit heat scheme
conserves its discrete sum even at an unstable time step. A mass check can
pass while oscillating errors grow. Conversely, a nonconserving boundary
condition can legitimately change total mass.

Record a normed error against a known solution on successive grids, estimate
the ratio, and verify that boundary and time-step conventions refine together.
A ratio near the expected value supports the observed regime. Roundoff,
singularities, or unresolved scales can prevent that regime from appearing.
Chapter~\ref{learn:probability} supplies uncertainty analysis when random
samples enter the computation; that uncertainty differs from deterministic
mesh error.

\section{Exercises with full solutions}
\paragraph{1. Choose an integration resolution.}
For the quadratic trapezoidal example, find the smallest positive integer
$n$ guaranteeing error at most $10^{-4}$.

\paragraph{Solution.}
The exact error is $1/(6n^2)$. The condition requires
$n^2\geq10000/6$, hence $n\geq41$. At $n=40$ the error is
$1/9600>10^{-4}$, whereas at $n=41$ it is $1/10086<10^{-4}$.
Thus 41 pieces suffice and are minimal for this specified rule and integrand.

\paragraph{2. Bound an iteration count.}
In the two-variable example, how many updates guarantee Euclidean error
below $0.01$ from $(0,0)$?

\paragraph{Solution.}
The initial error norm is $\sqrt2$ and the bound is $2^{-m}\sqrt2$.
At $m=7$ this is approximately $0.01105$; at $m=8$ it is approximately
$0.00552$. Eight updates suffice. The initial error lies in an eigenvector
direction of magnitude-one-half eigenvalue, so the bound is exact here.

\paragraph{3. Interpret a refinement table.}
Errors at spacings $h,h/2,h/4$ are $0.04,0.01,0.0025$. Estimate the
observed order and state the conclusion's boundary.

\paragraph{Solution.}
For $E(h)\approx Ch^p$, the ratio $E(h)/E(h/2)\approx2^p$.
Both measured ratios are four, giving $p=\log_2 4=2$.
The three measurements support second-order behavior on those spacings for
the specified test. They alone establish neither every-grid convergence
nor correctness of the physical model.
